\section{Requirements}
\label{sec:mercury-requirements}

In this section we describe the requirements of Mercury.

The developer specifies the interface to a module using MXDR (described later). The following
requirements apply:

\begin{itemize}
\item \textbf{Mercury.IDL}. Mercury shall take as input a file describing the messages of a
CubedOS module in \textit{Modified XDR} (MXDR) format.

\item \textbf{Mercury.Output}. Mercury shall produce an Ada package that defines the API to the
module. The package shall provide subprograms for encoding messages, decoding messages, and
classifying messages.

\item \textbf{Mercury.Provable}. Mercury shall generate \SPARK\ code that can be proved free of
runtime error without any further intervention by the user.
\end{itemize}

\subsection{Competitive Analysis}

The Core Flight Executive (cFE \cite{cFE}) passes space packets between modules with no special
internal encoding/decoding procedures for separate `proprietary messages.'

Most primary competitors to Mercury, such as ONC RPC, Ice, Java RMI, and WSDL, take a formal
language (generally a data representation language) and generate \newterm{stubs} that package
message parameters into message bodies on the client side,  and generate \newterm{skeletons}
that decompose messages on the server side back into the original parameters.

\begin{itemize}
\item \textbf{ONC RPC} is similar to Mercury in that it uses an IDL that is compiled into XDR
encoding stubs and XDR decoding skeletons. However, ONC RPC operates in a more complex
environment and contains many more features that are not relevant in the case of CubedOS.

\item \textbf{Ice} uses a proprietary language called Slice which can then be compiled to a
variety of languages including C++, Java, C\#, Python, Ruby, and PHP. It uses traditional
network protocols like TCP/IP and UDP/IP as the underlying transport.

\item \textbf{Java RMI} is like Ice, but only uses Java and the JVM for two-way data transfers.

\item \textbf{WSDL} is like Java RMI, but uses WSDL (Web Services Description Language).
\end{itemize}

\subsection{Base Requirements}

The following requirements apply:

\begin{itemize}
\item \textbf{Mercury.Robustness}. Mercury must be robust enough to handle all conceivable forms
of error in the MXDL file and output a clear and correct message describing each error.

\item \textbf{Mercury.Simplicity}. Mercury must be simple enough that a reasonable (and even
novice) user could successfully insert encoding and decoding subprograms into their CubedOS
project with minor to no difficulty.

\item \textbf{Mercury.XDR}. Mercury shall implement the whole of the XDR specification
\cite{rfc-4506}.
\end{itemize}

\subsection{XDR Extensions}

We refer you to the XDR specification \cite{rfc-4506} for the full details of the XDR encoding
that Mercury shall support in accordance with the \textbf{Mercury.XDR} requirement. In this
section, we focus on the extensions to XDR that Mercury shall support.

\todo{Describe the following: constrained typedefs, message structs, message aspects (e.g., message invariants).}

\subsection{Future Work}

In the future, Mercury could support other encoding schemes besides XDR (for example, ASN.1, to
name just one possibility). Also, a future version of Mercury could generate API packages that
use standard network protocols, such as TCP/IP (or some other), encoding messages into, for
example, a UDP datagram.
